
\subsubsection{2.1 Deduplication Effects on Language Model Performance}

Lee et al. (2022) provided the first systematic study showing that deduplicating training data via MinHash and exact substring methods reduces memorization and improves average model performance. This analysis was conducted at GPT-2 scale and reported aggregate averages across benchmarks. The per-benchmark distribution of effects --- whether contamination-proportional or uniform --- was not examined.

Muennighoff et al. (2023) analyzed data repetition effects on model training dynamics, finding diminishing returns at high repetition rates. This work focused on training dynamics rather than the contamination-performance correlation that is the present study's focus.

Biderman et al. (2023) introduced the Pythia suite and reported Pile versus dedup-Pile benchmark comparisons. That comparison used step-matching rather than token-count matching and did not perform contamination-accuracy correlation analysis. The present work extends Biderman et al. (2023) with a mechanistic analysis and methodological correction.

\subsubsection{2.2 Benchmark Contamination Detection and Measurement}

Brown et al. (2020) first characterized training data contamination in GPT-3 using n-gram overlap analysis and acknowledged that web-crawled training corpora may contain benchmark test content. The GPT-4 Technical Report adopted 13-gram overlap as a standard contamination measurement methodology.

Shi et al. (2023) introduced Min-k\% Probability (min-k\%), a likelihood-ratio-based method for detecting pretraining data membership. Min-k\% computes the minimum token probabilities over k\% of tokens in a sequence, providing a membership inference signal validated on multiple benchmarks. The present work applies min-k\% to the specific Pile/dedup-Pile distinction and finds an unexpected direction reversal at Pythia-1B scale.

Carlini et al. (2021) demonstrated that verbatim memorization in language models scales with model size. This scaling relationship is relevant to the present study's H-M2 direction reversal: the memorization signal may require model scales above 1B to manifest in the predicted direction.

Golchin and Surdeanu (2023) proposed data contamination quiz methods for instruction-tuned models. That work focuses on post-hoc contamination detection rather than characterizing performance effects as a function of contamination level.

\subsubsection{2.3 The Pile, dedup-Pile, and Pythia Infrastructure}

The Pile \citep{Gao2020Pile} is an 825GB English text corpus constructed from 22 diverse sources. Its deduplication variant, dedup-Pile, applies exact substring deduplication, reducing the corpus by approximately 15\% in token count. The Pythia suite \citep{Biderman2023Pythia} trained model families on both corpora under identical conditions, providing 154 intermediate checkpoints per model size (70M to 12B parameters). This checkpoint granularity enables the token-count matching methodology.

\subsubsection{2.4 Volume Effects and Scaling Laws}

Hoffmann et al. (2022) established compute-optimal scaling laws showing that model training performance depends on the total token count, not training steps alone. This provides the theoretical basis for the token-count matching methodology: at equal training steps, Pile and dedup-Pile models have processed different numbers of tokens, introducing a volume confound that step-matching fails to eliminate.

